% Options for packages loaded elsewhere
% Options for packages loaded elsewhere
\PassOptionsToPackage{unicode}{hyperref}
\PassOptionsToPackage{hyphens}{url}
%
\documentclass[
  11pt,
  ignorenonframetext,
  aspectratio=169,
]{beamer}
\newif\ifbibliography
\usepackage{pgfpages}
\setbeamertemplate{caption}[numbered]
\setbeamertemplate{caption label separator}{: }
\setbeamercolor{caption name}{fg=normal text.fg}
\beamertemplatenavigationsymbolsempty
% remove section numbering
\setbeamertemplate{part page}{
  \centering
  \begin{beamercolorbox}[sep=16pt,center]{part title}
    \usebeamerfont{part title}\insertpart\par
  \end{beamercolorbox}
}
\setbeamertemplate{section page}{
  \centering
  \begin{beamercolorbox}[sep=12pt,center]{section title}
    \usebeamerfont{section title}\insertsection\par
  \end{beamercolorbox}
}
\setbeamertemplate{subsection page}{
  \centering
  \begin{beamercolorbox}[sep=8pt,center]{subsection title}
    \usebeamerfont{subsection title}\insertsubsection\par
  \end{beamercolorbox}
}
% Prevent slide breaks in the middle of a paragraph
\widowpenalties 1 10000
\raggedbottom
\AtBeginPart{
  \frame{\partpage}
}
\AtBeginSection{
  \ifbibliography
  \else
    \frame{\sectionpage}
  \fi
}
\AtBeginSubsection{
  \frame{\subsectionpage}
}
\usepackage{iftex}
\ifPDFTeX
  \usepackage[T1]{fontenc}
  \usepackage[utf8]{inputenc}
  \usepackage{textcomp} % provide euro and other symbols
\else % if luatex or xetex
  \usepackage{unicode-math} % this also loads fontspec
  \defaultfontfeatures{Scale=MatchLowercase}
  \defaultfontfeatures[\rmfamily]{Ligatures=TeX,Scale=1}
\fi
\usepackage{lmodern}

\ifPDFTeX\else
  % xetex/luatex font selection
\fi
% Use upquote if available, for straight quotes in verbatim environments
\IfFileExists{upquote.sty}{\usepackage{upquote}}{}
\IfFileExists{microtype.sty}{% use microtype if available
  \usepackage[]{microtype}
  \UseMicrotypeSet[protrusion]{basicmath} % disable protrusion for tt fonts
}{}
\makeatletter
\@ifundefined{KOMAClassName}{% if non-KOMA class
  \IfFileExists{parskip.sty}{%
    \usepackage{parskip}
  }{% else
    \setlength{\parindent}{0pt}
    \setlength{\parskip}{6pt plus 2pt minus 1pt}}
}{% if KOMA class
  \KOMAoptions{parskip=half}}
\makeatother


\usepackage{longtable,booktabs,array}
\newcounter{none} % for unnumbered tables
\usepackage{calc} % for calculating minipage widths
\usepackage{caption}
% Make caption package work with longtable
\makeatletter
\def\fnum@table{\tablename~\thetable}
\makeatother
\usepackage{graphicx}
\makeatletter
\newsavebox\pandoc@box
\newcommand*\pandocbounded[1]{% scales image to fit in text height/width
  \sbox\pandoc@box{#1}%
  \Gscale@div\@tempa{\textheight}{\dimexpr\ht\pandoc@box+\dp\pandoc@box\relax}%
  \Gscale@div\@tempb{\linewidth}{\wd\pandoc@box}%
  \ifdim\@tempb\p@<\@tempa\p@\let\@tempa\@tempb\fi% select the smaller of both
  \ifdim\@tempa\p@<\p@\scalebox{\@tempa}{\usebox\pandoc@box}%
  \else\usebox{\pandoc@box}%
  \fi%
}
% Set default figure placement to htbp
\def\fps@figure{htbp}
\makeatother





\setlength{\emergencystretch}{3em} % prevent overfull lines

\providecommand{\tightlist}{%
  \setlength{\itemsep}{0pt}\setlength{\parskip}{0pt}}



 


\setbeamertemplate{navigation symbols}{}
\setbeamertemplate{footline}[frame number]
\setbeamertemplate{itemize items}[circle]
\setbeamertemplate{itemize subitem}[triangle]
\setbeamercolor{frametitle}{fg=black}
\setbeamerfont{frametitle}{series=\bfseries}
\usepackage{amsmath}
\usepackage{booktabs}
\usepackage{array}
\newcolumntype{R}[1]{>{\raggedright\arraybackslash}p{#1}}
\usepackage{fancyvrb}
\fvset{fontsize=\scriptsize}
\makeatletter
\def\verbatim@font{\scriptsize\ttfamily}
\makeatother
\makeatletter
\@ifpackageloaded{caption}{}{\usepackage{caption}}
\AtBeginDocument{%
\ifdefined\contentsname
  \renewcommand*\contentsname{Table of contents}
\else
  \newcommand\contentsname{Table of contents}
\fi
\ifdefined\listfigurename
  \renewcommand*\listfigurename{List of Figures}
\else
  \newcommand\listfigurename{List of Figures}
\fi
\ifdefined\listtablename
  \renewcommand*\listtablename{List of Tables}
\else
  \newcommand\listtablename{List of Tables}
\fi
\ifdefined\figurename
  \renewcommand*\figurename{Figure}
\else
  \newcommand\figurename{Figure}
\fi
\ifdefined\tablename
  \renewcommand*\tablename{Table}
\else
  \newcommand\tablename{Table}
\fi
}
\@ifpackageloaded{float}{}{\usepackage{float}}
\floatstyle{ruled}
\@ifundefined{c@chapter}{\newfloat{codelisting}{h}{lop}}{\newfloat{codelisting}{h}{lop}[chapter]}
\floatname{codelisting}{Listing}
\newcommand*\listoflistings{\listof{codelisting}{List of Listings}}
\makeatother
\makeatletter
\makeatother
\makeatletter
\@ifpackageloaded{caption}{}{\usepackage{caption}}
\@ifpackageloaded{subcaption}{}{\usepackage{subcaption}}
\makeatother

\usepackage{bookmark}
\IfFileExists{xurl.sty}{\usepackage{xurl}}{} % add URL line breaks if available
\urlstyle{same}
\hypersetup{
  pdftitle={Dynamic Panel Models},
  pdfauthor={Professor Ji-Woong Chung},
  hidelinks,
  pdfcreator={LaTeX via pandoc}}


\title{\texorpdfstring{Dynamic Panel Models}{Dynamic Panel Models}}
\subtitle{\texorpdfstring{BUSS975: Causal Inference in Financial
Research}{BUSS975: Causal Inference in Financial Research}}
\author{Professor Ji-Woong Chung}
\date{}
\institute{Korea University}

\begin{document}
\frame{\titlepage}


\begin{frame}{Where this lecture sits}
\protect\phantomsection\label{where-this-lecture-sits}
\begin{itemize}
\tightlist
\item
  Fixed effects removes everything permanent about a firm. The price:
  \emph{strict exogeneity}, no regressor may respond to past shocks.
\item
  The regressor that responds to past shocks by construction: the
  outcome's own lag. Leverage, payout, cash, investment and bank capital
  all depend on their own past.
\end{itemize}
\end{frame}

\section{1. Why the Lagged Outcome Is
Endogenous}\label{why-the-lagged-outcome-is-endogenous}

\begin{frame}{The model}
\protect\phantomsection\label{the-model}
\[
Y_{it} \;=\; \alpha\,Y_{i,t-1} + X_{it}'\beta + v_{it}, \qquad v_{it} = \mu_i + \varepsilon_{it}
\]

\begin{itemize}
\tightlist
\item
  \textbf{\(\alpha\): state dependence.} In partial-adjustment models
  (leverage targets, dividends, cash, bank capital), \(1-\alpha\) is the
  speed of adjustment.
\item
  \textbf{\(\mu_i\): the firm being itself.} Permanent; in leverage it
  explains most of the variation (Lemmon, Roberts and Zender 2008).
\item
  \textbf{\(\varepsilon_{it}\): news.} Mean zero, serially uncorrelated,
  unpredictable from anything known the year before.
\end{itemize}

\vspace{0.3em}

\begin{itemize}
\tightlist
\item
  \(Y_{i,t-1}\) is built from \(\mu_i\) at every past date, so pooled
  OLS, which leaves \(\mu_i\) in the error, is biased \textbf{up}. Last
  week's remedy: transform \(\mu_i\) away, by demeaning or by
  differencing. Does that settle it?
\end{itemize}
\end{frame}

\begin{frame}{Within estimator: demeaning borrows from the future}
\protect\phantomsection\label{within-estimator-demeaning-borrows-from-the-future}
Observations \(t = 1, \dots, T\), so \(Y_{i0}\) is the first lag.
Subtract each firm's time average from every variable: \[
\begin{gathered}
(Y_{it} - \bar Y_i) \;=\; \alpha\,(Y_{i,t-1} - \bar Y_{i,-1}) + (\varepsilon_{it} - \bar\varepsilon_i), \\[0.3em]
\bar Y_i = \tfrac{1}{T}\textstyle\sum_{s=1}^{T} Y_{is},\qquad
\bar Y_{i,-1} = \tfrac{1}{T}\textstyle\sum_{s=1}^{T} Y_{i,s-1},\qquad
\bar\varepsilon_i = \tfrac{1}{T}\textstyle\sum_{s=1}^{T} \varepsilon_{is}.
\end{gathered}
\]

\begin{itemize}
\tightlist
\item
  \(\mu_i\) is gone. But for \(t \ge 2\), \(\bar\varepsilon_i\) contains
  \(\tfrac{1}{T}\varepsilon_{i,t-1}\), so the demeaned error contains
  \(-\tfrac{1}{T}\varepsilon_{i,t-1}\), while the regressor
  \(Y_{i,t-1}\) contains \(+\varepsilon_{i,t-1}\): \[
  \operatorname{Cov}\!\big(Y_{i,t-1},\; \varepsilon_{it} - \bar\varepsilon_i\big) \;=\; -\frac{\sigma^2_\varepsilon}{T} \;<\; 0 \;\Longrightarrow\; \hat\alpha_{\text{FE}} \text{ biased \textbf{down}}.
  \]
\item
  A good shock last year raised \(Y_{i,t-1}\) \emph{and} raised the mean
  subtracted from every error. The regressor is high exactly when the
  transformed error has been pushed down. Demeaning uses the firm's
  whole history, future years included.
\end{itemize}
\end{frame}

\begin{frame}{First differencing: the overlap shrinks to one date}
\protect\phantomsection\label{first-differencing-the-overlap-shrinks-to-one-date}
\[
\Delta Y_{it} \;=\; \alpha\,\Delta Y_{i,t-1} + \Delta\varepsilon_{it}
\]

\(\mu_i\) is gone again. Every remaining object is a bundle of dated
shocks; write down which dates are inside each:

\begin{center}\footnotesize
\begin{tabular}{llc}
\toprule
Object & shocks dated & shares a date with $\Delta\varepsilon_{it}$? \\
\midrule
error $\Delta\varepsilon_{it} = \varepsilon_{it} - \varepsilon_{i,t-1}$ & $t,\; t-1$ & \\
regressor $\Delta Y_{i,t-1} = Y_{i,t-1} - Y_{i,t-2}$ & $t-1,\; t-2,\; \dots$ & \textbf{yes}: $t-1$ \\
$Y_{i,t-2}$ & $t-2,\; t-3,\; \dots$ & no \\
$\Delta Y_{i,t-2}$ & $t-2,\; t-3,\; \dots$ & no \\
\bottomrule
\end{tabular}
\end{center}

\begin{itemize}
\tightlist
\item
  \(\operatorname{Cov}(\Delta Y_{i,t-1}, \Delta\varepsilon_{it}) = -\operatorname{Var}(\varepsilon_{i,t-1}) = -\sigma^2_\varepsilon\):
  a shock at \(t-1\) raises the regressor and lowers the error. Biased
  \textbf{down}; more years do not help.
\end{itemize}

\begin{itemize}
\tightlist
\item
  Not the fix, but \textbf{where the fix starts}: \(\mu_i\) is gone and
  only a \emph{one-date} overlap remains. The last two rows share no
  date with the error.
\end{itemize}
\end{frame}

\section{2. Internal Instruments}\label{internal-instruments}

\begin{frame}{The assumption we can afford: sequential exogeneity}
\protect\phantomsection\label{the-assumption-we-can-afford-sequential-exogeneity}
\begin{itemize}
\item
  \textbf{Strict exogeneity}: shocks unrelated to every regressor at
  \emph{all} dates, past and future, \[
  \mathbb{E}\big[\varepsilon_{it} \,\big|\, X_{i1}, \dots, X_{iT}, \mu_i\big] = 0 \quad \text{(and likewise for } Y_{i,t-1}, \dots, Y_{i,T-1}).
  \] Fails by construction: \(\varepsilon_{it}\) moves \(Y_{it}\),
  tomorrow's regressor.
\item
  \textbf{Sequential exogeneity}: shocks unrelated to the \emph{past} of
  \(Y\) and the \emph{past and present} of \(X\), \[
  \mathbb{E}\big[\varepsilon_{it} \,\big|\, Y_{i,t-1}, \dots, Y_{i1},\; X_{it}, \dots, X_{i1},\; \mu_i\big] = 0.
  \] \textbf{\(\varepsilon_{it}\) is news}: it may change everything
  \emph{after} it (feedback allowed) but cannot be forecast from what
  was known \emph{before} it, that is, \(\varepsilon_{it}\) is serially
  uncorrelated.
\item
  Everything dated \(t-2\) or earlier was known before the news in
  \(\Delta\varepsilon_{it}\) arrived: \[
  \mathbb{E}\big[\,Y_{i,t-s}\,\Delta\varepsilon_{it}\,\big] = 0, \qquad s \ge 2.
  \] These \textbf{moment conditions} are the basis of what follows.
\end{itemize}
\end{frame}

\begin{frame}{Anderson and Hsiao (1981): the first fix}
\protect\phantomsection\label{anderson-and-hsiao-1981-the-first-fix}
Difference away \(\mu_i\), then instrument the one bad regressor,
\(\Delta Y_{i,t-1}\), by \textbf{2SLS}: \[
\Delta Y_{it} = \alpha\,\Delta Y_{i,t-1} + \Delta X_{it}'\beta + \Delta\varepsilon_{it}.
\]

\begin{center}\footnotesize
\begin{tabular}{R{2.6cm}R{4.3cm}R{4.6cm}l}
\toprule
instrument $z_{it}$ & relevant because & valid because & equations \\
\midrule
level $Y_{i,t-2}$ & it is part of $\Delta Y_{i,t-1} = Y_{i,t-1} - Y_{i,t-2}$ & dated $t-2$: $\mathbb{E}[Y_{i,t-2}\Delta\varepsilon_{it}] = 0$ & $t \ge 2$ \\
difference $\Delta Y_{i,t-2}$ & past changes predict recent changes & ends at $t-2$: $\mathbb{E}[\Delta Y_{i,t-2}\Delta\varepsilon_{it}] = 0$ & $t \ge 3$ \\
\bottomrule
\end{tabular}
\end{center}

\begin{itemize}
\tightlist
\item
  Apply 2SLS with one of the two as the instrument: just-identified,
  consistent as \(N \to \infty\) with \(T\) fixed.
\item
  \textbf{Its limits.} One instrument per equation, when at \(t = 5\)
  the levels \(Y_{i0}, \dots, Y_{i3}\) are all valid; and plain 2SLS
  ignores that \(\Delta\varepsilon_{it}\) and
  \(\Delta\varepsilon_{i,t-1}\) share a shock. Consistent but imprecise,
  which is why GMM replaced it.
\end{itemize}
\end{frame}

\begin{frame}{Difference GMM (Arellano and Bond 1991)}
\protect\phantomsection\label{difference-gmm-arellano-and-bond-1991}
\textbf{Step 1.} Difference away \(\mu_i\):
\(\;\Delta Y_{it} = \alpha\,\Delta Y_{i,t-1} + \Delta X_{it}'\beta + \Delta\varepsilon_{it}\).

\textbf{Step 2.} Instrument \(\Delta Y_{i,t-1}\) with lagged
\emph{levels}. Anderson--Hsiao used \(Y_{i,t-2}\) alone; every earlier
\(Y\) is just as \textbf{valid} (dated before any shock in
\(\Delta\varepsilon_{it}\)) and still \textbf{relevant}, so use them
all:

\begin{columns}[T]
\begin{column}{0.42\linewidth}
\begin{center}\small
\begin{tabular}{ll}
\toprule
equation at & instruments for $\Delta Y_{i,t-1}$ \\
\midrule
$t = 2$ & $Y_{i0}$ \\
$t = 3$ & $Y_{i0},\, Y_{i1}$ \\
$t = 4$ & $Y_{i0},\, Y_{i1},\, Y_{i2}$ \\
$t$ & $Y_{i0}, \dots, Y_{i,t-2}$ \\
\bottomrule
\end{tabular}
\end{center}
\end{column}

\begin{column}{0.56\linewidth}
\[
Z_i = \begin{pmatrix}
Y_{i0} & 0 & 0 & 0 & 0 & 0 \\
0 & Y_{i0} & Y_{i1} & 0 & 0 & 0 \\
0 & 0 & 0 & Y_{i0} & Y_{i1} & Y_{i2}
\end{pmatrix}
\quad
\begin{matrix} t = 2 \\ t = 3 \\ t = 4 \end{matrix}
\]
\end{column}
\end{columns}

\begin{itemize}
\tightlist
\item
  Anderson--Hsiao had one condition per equation; now there are many
  more conditions than parameters. What do you \emph{do} with them?
\end{itemize}
\end{frame}

\begin{frame}{What GMM does with the moment conditions}
\protect\phantomsection\label{what-gmm-does-with-the-moment-conditions}
Stack the instruments and the differenced residuals of firm \(i\) and
average over firms: \[
\bar g(\alpha) \;=\; \frac{1}{N}\sum_{i=1}^N Z_i'\,\Delta\varepsilon_i(\alpha), \qquad
\hat\alpha \;=\; \arg\min_\alpha\; J(\alpha) = N\,\bar g(\alpha)'\,W\,\bar g(\alpha).
\]

\begin{itemize}
\tightlist
\item
  Each entry of \(\bar g(\alpha)\) is the sample covariance between one
  instrument and the residual.
\item
  \textbf{The Overidentification Problem:} Validity says the
  \emph{population} covariances are zero. But with more instruments than
  parameters, no single \(\alpha\) can zero out all the \emph{sample}
  covariances at once.
\item
  \textbf{The GMM Solution:} Pick the \(\alpha\) that makes them
  \emph{jointly as close to zero as possible}.
\item
  Each instrument on its own gives a ``just-identified'' estimate. GMM
  is simply a weighted average of these estimates. The matrix \(W\)
  downweights the noisy, high-variance instruments (1-step vs 2-step
  GMM).
\end{itemize}
\end{frame}

\begin{frame}{Weak instruments: when a level cannot predict a change}
\protect\phantomsection\label{weak-instruments-when-a-level-cannot-predict-a-change}
\setlength{\abovedisplayskip}{3pt}\setlength{\belowdisplayskip}{3pt}

Difference GMM asks a \emph{level} \(Y_{i,t-2}\) to predict a
\emph{change} \(\Delta Y_{i,t-1}\). But if

\begin{itemize}
\tightlist
\item
  \textbf{\(\alpha \to 1\)}: little mean reversion; where a firm
  \emph{is} says nothing about where it is \emph{going}.
\item
  \textbf{\(\sigma^2_\mu\) large}: firms differ \emph{permanently}, and
  permanent differences predict no change.
\end{itemize}

So
\(\operatorname{Cov}(\Delta Y_{i,t-1}, Y_{i,t-2}) = -\sigma^2_\varepsilon/(1+\alpha)\)
stays bounded while
\(\operatorname{Var}(Y_{i,t-2}) = \sigma^2_\mu/(1-\alpha)^2 + \sigma^2_\varepsilon/(1-\alpha^2)\)
explodes; the first-stage correlation collapses:

\begin{center}\footnotesize
\begin{tabular}{lccccc}
\toprule
$\sigma^2_\mu / \sigma^2_\varepsilon$ & $\alpha = 0.5$ & $0.8$ & $0.9$ & $0.95$ & $0.99$ \\
\midrule
0 (no firm effect) & $-0.50$ & $-0.32$ & $-0.22$ & $-0.16$ & $-0.07$ \\
1 & $-0.25$ & $-0.10$ & $-0.05$ & $-0.02$ & $-0.00$ \\
4 & $-0.14$ & $-0.05$ & $-0.03$ & $-0.01$ & $-0.00$ \\
\bottomrule
\end{tabular}
\end{center}

\begin{itemize}
\tightlist
\item
  Bottom-right is finance: large standard errors, bias toward the within
  estimate.
\item
  \textbf{Blundell and Bond's answer.} A change \textbf{always predicts
  a level}: swap roles.
\end{itemize}
\end{frame}

\begin{frame}{System GMM (Blundell and Bond 1998)}
\protect\phantomsection\label{system-gmm-blundell-and-bond-1998}
Keep the differenced equation and \textbf{bring back the levels
equation}:

\begin{center}\footnotesize
\begin{tabular}{lll}
\toprule
equation & regressor & instruments \\
\midrule
differences: $\Delta Y_{it} = \alpha\Delta Y_{i,t-1} + \Delta\varepsilon_{it}$ & $\Delta Y_{i,t-1}$ & lagged \emph{levels} $Y_{i,t-2}, Y_{i,t-3}, \dots$ \\
levels: $Y_{it} = \alpha Y_{i,t-1} + \mu_i + \varepsilon_{it}$ & $Y_{i,t-1}$ & lagged \emph{differences} $\Delta Y_{i,t-1}$ \\
\bottomrule
\end{tabular}
\end{center}

\begin{itemize}
\tightlist
\item
  \textbf{Relevance does not fade.} The regressor contains the
  instrument, \(Y_{i,t-1} = Y_{i,t-2} + \Delta Y_{i,t-1}\): a firm that
  just moved up is now higher, whatever \(\alpha\) and \(\sigma^2_\mu\);
  no mean reversion needed.
\item
  \textbf{Requires a new assumption:} The levels error still contains
  the firm fixed effect \(\mu_i\). We need
  \(\operatorname{Cov}(\Delta Y_{i,t-1}, \mu_i) = 0\).

  \begin{itemize}
  \tightlist
  \item
    A firm's recent change says nothing about its permanent level.
  \item
    \emph{It fails:} When firms are still converging to their steady
    state (e.g., IPOs, post-crisis panels, post-deregulation).
  \end{itemize}
\item
  \textbf{The trade:} Difference GMM needs \emph{news} (low
  persistence). System GMM needs \emph{settled firms} (steady state).
\end{itemize}
\end{frame}

\begin{frame}{The simulation: five estimators in three worlds}
\protect\phantomsection\label{the-simulation-five-estimators-in-three-worlds}
\begin{center}
\includegraphics[width=1\linewidth,height=\textheight,keepaspectratio]{06b_dynamic_panel_models_files/figure-beamer/worlds-1.pdf}
\end{center}

\scriptsize \(500\) firms, \(T = 6\), \(100\) replications. Dots:
medians; bars: 10th--90th percentiles; dashed: truth. OLS biased up, FE
and FD down. \normalsize

\begin{itemize}
\tightlist
\item
  \textbf{\(\alpha = 0.5\).} Both GMM estimators sit on the truth
  (difference \(0.49\), system \(0.51\)).
\item
  \textbf{\(\alpha = 0.9\).} Difference GMM wide and low (\(0.77\)):
  weak instruments. System GMM tight and right (\(0.92\)).
\item
  Firms start away from steady state \(\mu_i\). Difference GMM recovers
  (\(0.89\)); system GMM lands at \(1.15\), above even OLS
\item
  Report OLS and FE beside GMM; outside their range is a warning.
\end{itemize}
\end{frame}

\begin{frame}{Summary}
\protect\phantomsection\label{summary}
\begin{enumerate}
\tightlist
\item
  \textbf{Removing \(\mu_i\) makes the lagged outcome endogenous.}
\item
  \textbf{Sequential exogeneity replaces strict exogeneity}: every \(Y\)
  dated \(t-2\) or earlier is orthogonal to \(\Delta\varepsilon_{it}\).
  Anderson--Hsiao uses one such lag by 2SLS; difference GMM uses them
  all and makes the sample moments as close to zero as \(W\) allows.
\item
  \textbf{The instruments are weak when the series is persistent or firm
  effects are large}. System GMM instruments the level with the change
  instead, at the price of mean stationarity, an assumption you must
  defend.
\item
  \textbf{In practice}: report OLS and FE beside the GMM column. OLS is
  biased up and FE down, so a GMM estimate outside their range signals a
  failed assumption. The diagnostics that test the assumptions directly,
  AR(2), Hansen and the instrument count, are in Roodman (2009) and
  Flannery and Hankins (2013).
\end{enumerate}
\end{frame}

\begin{frame}{References}
\protect\phantomsection\label{references}
\tiny

\begin{itemize}
\tightlist
\item
  Anderson, T. W. and C. Hsiao (1981). Estimation of Dynamic Models with
  Error Components. \emph{Journal of the American Statistical
  Association} 76, 598--606.
\item
  Arellano, M. and S. Bond (1991). Some Tests of Specification for Panel
  Data: Monte Carlo Evidence and an Application to Employment Equations.
  \emph{Review of Economic Studies} 58, 277--297.
\item
  Blundell, R. and S. Bond (1998). Initial Conditions and Moment
  Restrictions in Dynamic Panel Data Models. \emph{Journal of
  Econometrics} 87, 115--143.
\item
  Flannery, M. and K. Hankins (2013). Estimating Dynamic Panel Models in
  Corporate Finance. \emph{Journal of Corporate Finance} 19, 1--19.
\item
  Lemmon, M., M. Roberts and J. Zender (2008). Back to the Beginning:
  Persistence and the Cross-Section of Corporate Capital Structure.
  \emph{JF} 63, 1575--1608.
\item
  Nickell, S. (1981). Biases in Dynamic Models with Fixed Effects.
  \emph{Econometrica} 49, 1417--1426.
\item
  Roodman, D. (2009). How to Do xtabond2: An Introduction to Difference
  and System GMM in Stata. \emph{Stata Journal} 9, 86--136; and A Note
  on the Theme of Too Many Instruments. \emph{Oxford Bulletin of
  Economics and Statistics} 71, 135--158.
\end{itemize}
\end{frame}




\end{document}
